% Appendix H -- the superseded four-point rider sweep, moved out of section 6
% (2026-09-26). Kept rather than deleted because the paper's earlier overlap
% figure (78%) came from it, and a reader who has seen that number deserves to
% see what replaced it and why.

\section{The superseded rider sweep}
\label{app:ridercost}

Table~\ref{tab:ridercost} is the four-point sweep \S\ref{sec:eval:cost} sets
aside, kept here because the 78\% overlap figure it implies appeared in earlier
drafts and a reader who has met that number should see what replaced it.

Sweeping the number of host-tier requests on an otherwise identical engine and
fitting
$\mathrm{ITL}(r) = \mathrm{ITL}_0 + \text{fixed} + r\cdot c$
separates a fixed per-step cost of using the host path at all from a per-request
cost that fails to overlap compute. The two extrapolate to completely different
answers on a faster link, which is why a single operating point cannot settle
the question.

\begin{table}[t]
\centering
\small
\begin{tabular}{lrr}
\toprule
host-tier requests & ITL (ms) & tok/s \\
\midrule
0 & 49.99 & 57.5 \\
1 & 65.66 & 54.7 \\
2 & 65.93 & 55.8 \\
3 & 74.98 & 57.0 \\
\bottomrule
\end{tabular}
\caption{Qwen2.5-32B, 2048 context, 12 offered, 8 resident, A100/PCIe.
\emph{ITL} is the median inter-token latency; \emph{tok/s} is end-to-end decode
throughput over the whole 12-request run including the drain, \emph{not}
$B/\mathrm{ITL}$. \textbf{The arms do not hold the resident cohort fixed}: peak
concurrency runs 4, 8, 10, 11 down the rows, so $\Delta$ITL conflates a rider's
cost with additional \emph{resident} requests --- the confound
\S\ref{sec:eval:ablation} calls disqualifying. ``8 resident'' is the cap, not
the observed cohort. An affine fit over these four points gives $9.6$\,ms fixed
$+$ $4.7$\,ms per request, implying $78\%$ overlap; the direct measurement below
supersedes both that functional form and that number, and the fit is three
parameters' worth of structure read from four points.}
\label{tab:ridercost}
\end{table}
% src: paper/data/hide_efficiency.json @ f3df05bdb

\textbf{The throughput column is end-to-end over a finite workload, and that
matters.} It is total decode tokens divided by total decode wall time across all
12 requests, including the drain at the end, so it is not $B/\mathrm{ITL}$ and
does not reconcile with the ITL column: at $r{=}0$, $768/13.4\,\mathrm{s}$ gives
57.3 tok/s while $B/\mathrm{ITL}$ at the peak concurrency of 4 would give 80.
The two readings of these same four rows disagree about the sign. End to end,
throughput is flat within noise (57.5, 54.7, 55.8, 57.0). Read as
$B/\mathrm{ITL}$ at peak concurrency it \emph{rises} (80, 122, 152, 147),
because riders add concurrency faster than they lengthen the step.

Neither is decisive, and the reason is the experiment rather than the system: a
12-request closed workload never reaches steady state, so the end-to-end column
is dominated by the drain, while peak concurrency is an upper bound that no step
sustains. A third reading of a single two-request point is worse still --- it
cannot distinguish a steep slope from a large intercept, and read that way it
implies 89\% of the transfer fails to overlap, which the sweep contradicts. What
settles the cost is an experiment that holds the resident batch fixed and adds
riders on top, which is the measurement \S\ref{sec:eval:cost} keeps.

The fixed-versus-slope split is also unstable at three points: two runs of the
same sweep fitted $12.3 + 4.5$ and $9.6 + 4.7$. Because the fixed term is
precisely the part that a faster interconnect does \emph{not} shrink, that
split determines break-even on a coherent link, and three points cannot fix it.
Any interconnect comparison therefore needs the sweep run at ten host-tier
counts or more; the curve in \S\ref{sec:eval:cost} is the first step of that and still stops at
six.
